\section{Polynesian}

The Polynesian traditions are the cosmologies of the peoples spread across the islands of the central and southern Pacific. They share a common ancestral language and a common stock of gods, told in local forms across Hawaii, New Zealand, and the wider region. Two accounts stand out. The Hawaiian \textbf{Kumulipo}, a long chant of emergence from the deep darkness, and the Maori story of the sky father and the earth mother held together and then pulled apart by their children.

\subsection{The Creation Model}

In the first account the world begins in a deep and slimy darkness that is not empty but full of latent life. From that darkness the first living thing arises, a small coral animal that builds the reef. After it come, in ordered stages, the creatures of the sea, then the crawling and flying things of the land, then at last the gods and humankind. Creation is a slow climb out of the dark and toward the light, each new kind born from the one before it.

In the second account the sky and the earth are persons, a father and a mother, locked in a tight embrace. Their children are crushed in the narrow dark between them. The children take counsel. One after another they try to force the parents apart and fail, until the god of the forests braces his back against the earth and his feet against the sky and pushes them open. Light floods in and the world of the living begins. One angry child, the god of winds, will not accept the parting and storms against his brothers forever.

\subsection{Key Motifs}

The Kumulipo runs in sixteen sections called \textbf{wa}. The first seven belong to \textbf{po}, the dark spirit-realm, and the last nine to \textbf{ao}, the realm of light. The Maori frame the whole movement as three states.

\begin{longtable}{L{0.22\linewidth}
L{0.6\linewidth}}
\toprule
\textbf{State} & \textbf{Meaning} \\
\midrule
\textbf{Te Kore} & the void, the nothing that holds all potential \\
\textbf{Te Po} & the long night, the dark of forming \\
\textbf{Te Ao} & the world of light, the present living world \\
\bottomrule
\end{longtable}

The sky father is \textbf{Ranginui} and the earth mother is \textbf{Papatuanuku}. The forest god who parts them is \textbf{Tane}, and the storm god who resists is \textbf{Tawhirimatea}.

This tradition states the shared pattern as a passage from the one dark source out into the many forms of light. The going-out is the climb from the deep darkness, or the prying open of the joined sky and earth, and the return is held in the grief of the parted parents who still reach toward each other.
